\section{Related Work}
\label{sec:related}

We did not run a systematic literature search for this study; this section positions our work against the papers that motivated it.

\para{Chain-of-thought prompting.}
\citet{wei2022chain} show that few-shot exemplars containing intermediate reasoning steps improve performance on arithmetic, commonsense and symbolic tasks, and \citet{kojima2022large} show that a single zero-shot instruction to think step by step yields similar gains.
We use the zero-shot form, because it is what practitioners apply by default and it needs no task-specific exemplars.

\para{Where \CoT helps.}
\citet{sprague2024cot} conduct a large meta-analysis and new experiments and find that \CoT gains concentrate on math and symbolic reasoning, with little benefit elsewhere.
Their finding for other tasks is mainly an absence of gain.
Building on this, we ask the sharper question of whether \CoT is \emph{harmful} on recognition and recall tasks, and we use paired item-level outcomes to separate a null effect from offsetting harms and fixes.

\para{Where \CoT hurts.}
\citet{liu2024mind} show that \CoT can reduce performance on tasks where verbal deliberation makes humans worse, such as implicit pattern learning.
Unlike their task selection, our ``pattern matching'' category consists of standard text classification, which models already solve near ceiling.
Our null result on classification therefore does not contradict theirs; it shows that the harm they find does not extend to ordinary classification in the models we test.

\para{Scale.}
Most comparisons of \CoT across model sizes span different families, so scale is confounded with training data and recipe.
We hold the family fixed and vary only size, using the five smallest instruction-tuned sizes of Qwen2.5 \cite{qwen2024qwen25}.
